\section*{Item 3}\label{it:3}

\textbf{Enunciado}: Explique:

\medskip
\textbf{a)} a diferença entre o funcionamento do elemento magnético em conversores isolados como indutores acoplados e transformadores.

\medskip
\textbf{Resolução}: em conversores isolados, o elemento magnético atua como indutores acoplados. A energia é armazenada no campo magnético durante uma etapa de operação e posteriormente transferida para a saída, sendo necessário considerar a capacidade da indutância de magnetização. Dessa forma, o indutor acoplado opera principalmente como elemento de isolamento galvânico e adequação do nível de tensão.

\medskip
\textbf{b)} a diferença no aproveitamento da curva $B$--$H$ nos conversores Flyback, Forward e Full Bridge.

\medskip
\textbf{Resolução}: o aproveitamento da curva $B$--$H$ depende da forma como ocorre a magnetização e a desmagnetização do elemento magnético.

No conversor Flyback, o núcleo opera com excursão em apenas uma região da curva $B$--$H$. Dessa forma, o fluxo magnético varia entre valores positivos, sem utilizar toda a região disponível antes da saturação do núcleo, resultando em menor aproveitamento do material magnético.

No conversor Forward, após o desligamento da chave, é necessário realizar a desmagnetização total do núcleo, permitindo que o fluxo retorne a zero. Assim, a excursão de fluxo é maior que no Flyback, porém ainda ocorre predominantemente em um único sentido da curva $B$--$H$.

No conversor Full Bridge, as chaves permitem aplicar tensão alternada no primário do indutor acoplado, fazendo com que o fluxo magnético varie em ambos os sentidos. Dessa forma, o núcleo utiliza os dois quadrantes da curva $B$--$H$, proporcionando melhor aproveitamento do material magnético e permitindo maiores níveis de potência.

\medskip
\textbf{c)} em uma aplicação de carregamento rápido de veículos elétricos, qual dos conversores apresentados em aula é o mais indicado e por quê?

\medskip
\textbf{Resolução}: para aplicações de carregamento rápido de veículos elétricos, o conversor Full Bridge é o mais indicado entre as topologias apresentadas. Isso ocorre devido à seu melhor aproveitamento do núcleo magnético e menor estresse nos componentes quando comparado às topologias Flyback e Forward. Além diso, a operação em ambos os sentidos da curva $B$--$H$ reduz o volume do transformador para uma mesma potência processada.